\section{应用公司与模型公司的边界会变模糊}

上一章讲生成系统栈，本章讨论公司边界。王冠认为应用公司和模型公司的边界会变得模糊。原因是模型公司会向上做产品，应用公司会向下做数据、DSL、context 和局部模型优化。真正的差异不在公司标签，而在谁拥有最关键的数据飞轮和用户场景。

\begin{figure}[H]
\centering
\includegraphics[width=0.94\textwidth]{app-model-boundary.png}
\caption{应用公司与模型公司边界：基座模型与产品内生数据分别形成不同飞轮。自制概念图，依据 01:41:59--02:01:44 对谈内容整理。}
\end{figure}

\begin{knowledgebox}{读图：边界模糊不等于边界消失}
模型公司拥有基座模型、算力、通用能力和 API 分发；应用公司拥有场景、产品内生数据、工作流和用户价值。未来竞争的关键，是谁能把模型能力和场景数据接成闭环。
\end{knowledgebox}

\subsection{模型公司为什么会上探产品}

本节先看模型公司向应用层移动的原因。模型能力如果只停留在 API，会离用户反馈和收入上限太远；产品入口能把模型能力变成更高价值的商业关系。

模型公司如果只卖 API，很容易被价格和算力成本约束；上探产品可以拿到用户入口、真实反馈和更高收入。ChatGPT、Claude Code、Gemini 等都在证明模型公司会直接进入应用层。尤其是 Coding 这类场景，模型公司天然掌握数据和开发者入口，更容易自己做。

\subsection{应用公司为什么要下探模型和数据}

接下来从应用公司角度看同一个问题。如果模型公司会上探产品，应用公司就不能只停留在皮肤层，而要向下构建数据、方法和工作流。

应用公司如果只做 UI，也会被模型公司复制。它必须向下构建 DSL、context、领域数据、评测、反馈和局部优化。它不一定要训练通用大模型，但需要让自己的生成系统拥有别人拿不到的产品内生数据。

\begin{importantbox}{应用公司的护城河}
应用公司的护城河不只是界面，而是“场景中的数据飞轮”：用户意图、生产过程、失败样本、recipe、商业反馈和持续评测。
\end{importantbox}

\subsection{本章小结}

应用公司和模型公司的边界会模糊，但不会没有差异。模型公司从能力出发，应用公司从场景出发；谁能形成更强的数据飞轮，谁就能在生成系统时代获得更稳的位置。

